\chapter{Sheaves of Differentials}

\section{Kahler Differentials}

\begin{definition}
Let \(A\to B\) be a ring homomorphism and let \(M\) be a \(B\)-module.  An
\(A\)-\emph{derivation} from \(B\) to \(M\) is an \(A\)-linear map
\(d:B\to M\) such that
\[
    d(bb')=b\,d b'+b'\,d b
\]
for all \(b,b'\in B\).  The group of \(A\)-derivations is denoted
\(\Der_A(B,M)\).
\end{definition}

\begin{definition}
The module of \emph{Kahler differentials} \(\Omega_{B/A}\) is the \(B\)-module
generated by symbols \(db\), \(b\in B\), subject to the relations
\[
    d(b+b')=db+db',\qquad d(bb')=b\,db'+b'\,db,\qquad da=0
\]
for \(a\in A\).
\end{definition}

\begin{proposition}[Universal property]
For every \(B\)-module \(M\), composition with the universal derivation
\[
    d:B\longrightarrow \Omega_{B/A},\qquad b\longmapsto db,
\]
induces a natural isomorphism
\[
    \Hom_B(\Omega_{B/A},M)\simeq \Der_A(B,M).
\]
\end{proposition}

\begin{proof}
If \(\varphi:\Omega_{B/A}\to M\) is \(B\)-linear, then
\(\varphi\circ d\) satisfies the defining rules for an \(A\)-derivation because
the relations in \(\Omega_{B/A}\) impose additivity, the Leibniz rule, and
vanishing on \(A\).  Conversely, if \(\delta:B\to M\) is an \(A\)-derivation,
then the assignment \(db\mapsto\delta(b)\) respects exactly the defining
relations, so it extends uniquely to a \(B\)-linear map
\(\Omega_{B/A}\to M\).  These two constructions are inverse.
\end{proof}

\begin{proposition}
If \(B=A[x_1,\ldots,x_n]\), then \(\Omega_{B/A}\) is a free \(B\)-module with
basis \(dx_1,\ldots,dx_n\).
\end{proposition}

\begin{proof}
Let \(M\) be a \(B\)-module.  An \(A\)-derivation \(\delta:B\to M\) is uniquely
determined by the elements \(\delta(x_i)\), because the Leibniz rule determines
\(\delta\) on monomials:
\[
    \delta(x_1^{e_1}\cdots x_n^{e_n})
    =
    \sum_i e_i x_1^{e_1}\cdots x_i^{e_i-1}\cdots x_n^{e_n}\delta(x_i).
\]
Conversely, any choice of \(m_i\in M\) defines a derivation by the displayed
formula and linearity.  Hence
\[
    \Der_A(B,M)\simeq M^n\simeq
    \Hom_B\left(\bigoplus_{i=1}^n B\,dx_i,M\right)
\]
naturally in \(M\).  By the universal property, \(\Omega_{B/A}\) is the free
module with basis \(dx_i\).
\end{proof}

\begin{proposition}
Let \(B=A[x_1,\ldots,x_n]/I\).  Then there is an exact sequence
\[
    I/I^2\longrightarrow
    \bigoplus_{i=1}^n B\,dx_i
    \longrightarrow \Omega_{B/A}\longrightarrow 0,
\]
where the first map sends the class of \(f\in I\) to
\(\sum_i (\partial f/\partial x_i)\,dx_i\).
\end{proposition}

\begin{proof}
Let \(P=A[x_1,\ldots,x_n]\).  Any \(A\)-derivation \(B\to M\) is the same as an
\(A\)-derivation \(\delta:P\to M\) that vanishes on \(I\), where \(M\) is viewed
as a \(P\)-module through \(P\to B\).  Since \(\Omega_{P/A}\) is free with basis
\(dx_i\), such derivations are the \(B\)-linear maps
\(\Omega_{P/A}\tensor_P B\to M\) that kill the image of \(I/I^2\).  Therefore
they are the same as \(B\)-linear maps from the cokernel of
\[
    I/I^2\to\Omega_{P/A}\tensor_P B.
\]
By the universal property, this cokernel is \(\Omega_{B/A}\).
\end{proof}

\section{Differentials of Schemes}

\begin{definition}
Let \(f:X\to S\) be a morphism of schemes.  The sheaf of relative differentials
\(\Omega_{X/S}\) is the quasi-coherent \(\OO_X\)-module obtained by gluing the
modules \(\widetilde{\Omega_{B/A}}\) on affine opens
\(\Spec B\subset X\) mapping into affine opens \(\Spec A\subset S\).
\end{definition}

\begin{proposition}
The sheaf \(\Omega_{X/S}\) represents \(S\)-derivations: for every
\(\OO_X\)-module \(\MM\), there is a natural isomorphism
\[
    \Hom_{\OO_X}(\Omega_{X/S},\MM)
    \simeq
    \Der_S(\OO_X,\MM).
\]
\end{proposition}

\begin{proof}
On affine opens \(\Spec B\subset X\) over \(\Spec A\subset S\), this is the
universal property of \(\Omega_{B/A}\).  These affine identifications are
compatible with localization, because localization commutes with Kahler
differentials:
\[
    \Omega_{B_g/A}\simeq \Omega_{B/A}\tensor_B B_g.
\]
The local maps therefore glue, giving the asserted global isomorphism.
\end{proof}

\begin{proposition}[Transitivity sequence]
For morphisms \(X\xto{f}Y\xto{g}S\), there is a natural exact sequence
\[
    f^*\Omega_{Y/S}\longrightarrow \Omega_{X/S}
    \longrightarrow \Omega_{X/Y}\longrightarrow 0.
\]
\end{proposition}

\begin{proof}
The assertion is local on \(X\) and \(Y\).  Let \(A\to B\to C\) be the
corresponding ring maps.  The map
\[
    \Omega_{B/A}\tensor_B C\longrightarrow\Omega_{C/A}
\]
sends \(db\tensor 1\) to \(d(b)\) in \(C\), and the map
\(\Omega_{C/A}\to\Omega_{C/B}\) sends \(dc\) to \(dc\).  The latter is
surjective because the symbols \(dc\) generate \(\Omega_{C/B}\).  Its kernel is
generated by the differentials of elements of \(B\), which are precisely the
image of \(\Omega_{B/A}\tensor_B C\).  Sheafifying gives the exact sequence.
\end{proof}

\begin{proposition}[Conormal sequence]
Let \(i:Z\hookrightarrow X\) be a closed immersion over \(S\), with ideal sheaf
\(\II\subset\OO_X\).  Then there is an exact sequence on \(Z\)
\[
    \II/\II^2\longrightarrow i^*\Omega_{X/S}
    \longrightarrow \Omega_{Z/S}\longrightarrow 0.
\]
\end{proposition}

\begin{proof}
The statement is local on \(X\).  Write \(X=\Spec B\) and
\(Z=\Spec C\), with \(C=B/I\).  Applying the presentation result for
differentials to the quotient map \(B\to C\) gives
\[
    I/I^2\longrightarrow \Omega_{B/A}\tensor_B C
    \longrightarrow \Omega_{C/A}\longrightarrow0.
\]
This is exactly the displayed sequence after passing to associated sheaves.
\end{proof}

\begin{proposition}[Base change]
Given a cartesian square
\[
\begin{tikzcd}
X' \arrow[r,"g'"] \arrow[d,"f'"'] & X \arrow[d,"f"]\\
S' \arrow[r,"g"] & S ,
\end{tikzcd}
\]
there is a natural morphism
\[
    g'^*\Omega_{X/S}\longrightarrow \Omega_{X'/S'}.
\]
If the square is obtained from affine rings by \(B'=B\tensor_A A'\), this map
is an isomorphism.
\end{proposition}

\begin{proof}
On affine opens the square is represented by \(A\to B\) and \(A\to A'\), with
\(B'=B\tensor_A A'\).  The universal derivation \(B\to\Omega_{B/A}\), after
tensoring with \(B'\), gives an \(A'\)-derivation
\[
    B'\longrightarrow \Omega_{B/A}\tensor_B B',
\]
so by the universal property there is a map
\[
    \Omega_{B'/A'}\longrightarrow \Omega_{B/A}\tensor_B B'.
\]
Conversely, the derivation \(B\to B'\to\Omega_{B'/A'}\) is \(A\)-linear, hence
gives a \(B'\)-linear map
\[
    \Omega_{B/A}\tensor_B B'\longrightarrow \Omega_{B'/A'}.
\]
The two maps send \(db\tensor1\) and \(d(b\tensor a')\) to the corresponding
universal differentials and are inverse on generators.  Thus the base-change
map is an isomorphism in the affine tensor product case, and these local
isomorphisms glue.
\end{proof}

\section{Smooth Curves and the Canonical Sheaf}

\begin{definition}
Let \(k\) be a field.  A scheme \(X\) of finite type over \(k\) is \emph{smooth
of dimension \(d\)} if every point has an affine neighbourhood
\(\Spec B\) such that, locally at the point, \(B\) has a presentation
\[
    k[x_1,\ldots,x_n]/(f_1,\ldots,f_{n-d})
\]
for which the Jacobian matrix
\((\partial f_i/\partial x_j)\) has rank \(n-d\).  A \emph{smooth curve} over
\(k\) is a smooth \(k\)-scheme of dimension one.
\end{definition}

\begin{proposition}
If \(X\) is smooth over \(k\) of dimension \(d\), then \(\Omega_{X/k}\) is
locally free of rank \(d\).  In particular, if \(X\) is a smooth curve, then
\(\Omega_{X/k}\) is an invertible sheaf.
\end{proposition}

\begin{proof}
Work locally with a presentation
\[
    B=k[x_1,\ldots,x_n]/(f_1,\ldots,f_{n-d})
\]
for which the Jacobian matrix has rank \(n-d\).  The conormal presentation gives
\[
    B^{n-d}\longrightarrow B^n\longrightarrow\Omega_{B/k}\to0,
\]
where the first map is the Jacobian matrix.  After possibly shrinking to a
principal open, some \((n-d)\times(n-d)\) minor is invertible.  Elementary row
and column operations then identify the map \(B^{n-d}\to B^n\) with the
inclusion of the first \(n-d\) basis vectors.  The cokernel is free of rank
\(d\).  These local computations glue, so \(\Omega_{X/k}\) is locally free of
rank \(d\).
\end{proof}

\begin{definition}
For a smooth curve \(X\) over \(k\), the \emph{canonical sheaf} is
\[
    \omega_X=\Omega_{X/k}.
\]
A \emph{rational differential} is an element of
\(\Omega_{k(X)/k}\), where \(k(X)\) is the function field of an integral curve
\(X\).
\end{definition}

\begin{proposition}
Let \(X\) be a smooth integral curve over \(k\).  For every closed point
\(P\in X\), the local ring \(\OO_{X,P}\) is a discrete valuation ring.
\end{proposition}

\begin{proof}
Since \(X\) is smooth of dimension one, \(\OO_{X,P}\) is a noetherian regular
local ring of dimension one.  A one-dimensional noetherian regular local ring
has maximal ideal generated by one non-zero element \(t\).  Every non-zero
ideal \(I\) contains an element of minimal \(t\)-adic order, say \(u t^n\) with
unit \(u\), and then \(I=(t^n)\).  Thus every non-zero ideal is a power of the
maximal ideal, which is the defining ideal-theoretic property of a discrete
valuation ring.  The associated valuation is denoted \(\ord_P\).
\end{proof}

\section{Cotangent Spaces and Unramified Morphisms}

\begin{proposition}
Let \(X=\Spec B\) be an affine scheme over a field \(k\), and let \(x\in X\) be
a \(k\)-rational point with maximal ideal \(\mathfrak m\subset B\).  Then there
is a canonical isomorphism
\[
    \Omega_{B/k}\tensor_B k\simeq \mathfrak m/\mathfrak m^2.
\]
\end{proposition}

\begin{proof}
The universal derivation \(d:B\to\Omega_{B/k}\) sends \(\mathfrak m\) to
\(\Omega_{B/k}\tensor_B k\).  Since \(d(ab)=a\,db+b\,da\), the image of
\(\mathfrak m^2\) is zero after tensoring with \(k=B/\mathfrak m\).  Thus there
is a map \(\mathfrak m/\mathfrak m^2\to\Omega_{B/k}\tensor_B k\).  Conversely,
every \(k\)-derivation \(B\to k\) vanishes on \(\mathfrak m^2\), hence factors
through a \(k\)-linear map \(\mathfrak m/\mathfrak m^2\to k\).  By the universal
property of differentials, the duals of the two vector spaces represent the
same functor of \(k\)-vector spaces.  Therefore they are canonically
isomorphic.
\end{proof}

\begin{definition}
A morphism \(f:X\to Y\) locally of finite type is \emph{unramified at}
\(x\in X\) if
\[
    \Omega_{X/Y}\tensor_{\OO_X}\kappa(x)=0
\]
and \(\kappa(x)\) is a finite separable extension of \(\kappa(f(x))\).  It is
\emph{unramified} if it is unramified at every point.
\end{definition}

\begin{proposition}
Let \(A\to B\) be a finite type morphism and let
\[
    B=A[x_1,\ldots,x_n]/(f_1,\ldots,f_r).
\]
For \(\mathfrak q\in\Spec B\), the vector space
\(\Omega_{B/A}\tensor_B\kappa(\mathfrak q)\) is the cokernel of the Jacobian
matrix
\[
    \left(\frac{\partial f_i}{\partial x_j}\right)(\mathfrak q):
    \kappa(\mathfrak q)^r\to\kappa(\mathfrak q)^n.
\]
\end{proposition}

\begin{proof}
The conormal presentation gives
\[
    I/I^2\to \bigoplus_j B\,dx_j\to\Omega_{B/A}\to0.
\]
Tensoring with \(\kappa(\mathfrak q)\) is right exact.  The first map sends
\(f_i\) to \(\sum_j(\partial f_i/\partial x_j)dx_j\).  Thus the fiber of
\(\Omega_{B/A}\) at \(\mathfrak q\) is the cokernel of the displayed matrix.
\end{proof}

\section{Smoothness Criteria}

\begin{theorem}[Jacobian criterion]
Let \(k\) be a perfect field and let
\[
    X=\Spec k[x_1,\ldots,x_n]/(f_1,\ldots,f_r)
\]
be of finite type over \(k\).  A \(k\)-rational point \(x\in X\) is smooth of
dimension \(d\) if and only if the Jacobian matrix
\[
    \left(\frac{\partial f_i}{\partial x_j}(x)\right)
\]
has rank \(n-d\), where \(d=\dim_x X\).
\end{theorem}

\begin{proof}
The tangent space has dimension
\[
    n-\rank\left(\frac{\partial f_i}{\partial x_j}(x)\right)
\]
by the tangent-space computation and the presentation of differentials.  For a
finite type algebra over a perfect field, the local ring at \(x\) is regular
exactly when the dimension of the Zariski tangent space equals the Krull
dimension of the local ring.  Substituting \(d=\dim_x X\) gives the claim.
\end{proof}

\begin{theorem}[Differential criterion for smoothness]
Let \(f:X\to S\) be locally of finite presentation.  If \(f\) is smooth of
relative dimension \(d\), then \(\Omega_{X/S}\) is locally free of rank \(d\).
Conversely, if \(f\) is flat, the fibers are geometrically reduced, and
\(\Omega_{X/S}\) is locally free of rank \(d\) with the expected fiber
dimension, then \(f\) is smooth of relative dimension \(d\).
\end{theorem}

\begin{proof}
The first statement is local and follows from the Jacobian presentation: after
choosing local equations with a maximal minor invertible, the conormal sequence
identifies \(\Omega_{X/S}\) with a free module of rank \(d\).  For the
converse, flatness reduces smoothness to regularity of geometric fibers.  On
each fiber, the module of differentials has dimension \(d\) at every point, so
the tangent space has dimension \(d\).  The expected fiber dimension is also
\(d\).  The regularity criterion for finite type algebras over fields gives
regular fibers, hence smoothness.
\end{proof}

\begin{corollary}
An etale morphism \(f:X\to S\) locally of finite presentation satisfies
\[
    \Omega_{X/S}=0.
\]
Conversely, a flat unramified morphism locally of finite presentation is etale.
\end{corollary}

\begin{proof}
Etale means smooth of relative dimension zero, so the preceding theorem gives
\(\Omega_{X/S}\) locally free of rank zero, hence zero.  Conversely, if \(f\) is
flat and unramified, all fibers are finite separable field extensions locally
and have dimension zero; the differential criterion gives smoothness of
relative dimension zero.
\end{proof}

\section{Higher Differential Forms}

\begin{definition}
For a morphism \(f:X\to S\), define
\[
    \Omega^p_{X/S}=\bigwedge^p_{\OO_X}\Omega_{X/S}.
\]
If \(X\to S\) is smooth of relative dimension \(d\), the \emph{relative
canonical sheaf} is
\[
    \omega_{X/S}=\Omega^d_{X/S}.
\]
\end{definition}

\begin{proposition}
If \(X\to S\) is smooth of relative dimension \(d\), then
\(\Omega^p_{X/S}\) is locally free of rank \(\binom{d}{p}\), and
\(\omega_{X/S}\) is invertible.
\end{proposition}

\begin{proof}
Locally \(\Omega_{X/S}\simeq\OO_X^{\oplus d}\).  The \(p\)-th exterior power of
a free module of rank \(d\) is free with basis indexed by \(p\)-element subsets
of \(\{1,\ldots,d\}\), hence has rank \(\binom{d}{p}\).  For \(p=d\), the rank
is one, so \(\omega_{X/S}\) is invertible.
\end{proof}

\section{Ramification of Maps of Curves}

\begin{definition}
Let \(f:X\to Y\) be a non-constant finite morphism of smooth projective curves.
For a closed point \(P\in X\), with \(Q=f(P)\), the \emph{ramification index}
\(e_P\) is defined by
\[
    f^\#(\mathfrak m_Q)\OO_{X,P}=\mathfrak m_P^{e_P}.
\]
\end{definition}

\begin{proposition}
Let \(f:X\to Y\) be a finite separable morphism of smooth curves.  There is an
exact sequence
\[
    f^*\Omega_{Y/k}\longrightarrow \Omega_{X/k}
    \longrightarrow \Omega_{X/Y}\longrightarrow0,
\]
and \(\Omega_{X/Y}\) is a torsion sheaf supported at the ramification points.
\end{proposition}

\begin{proof}
The exact sequence is the transitivity sequence for
\(X\to Y\to\Spec k\).  Since \(f\) is separable, the induced extension of
function fields is separable, so
\[
    \Omega_{k(X)/k(Y)}=0.
\]
Thus the generic stalk of \(\Omega_{X/Y}\) is zero.  On a noetherian curve, a
coherent sheaf with zero generic stalk is a torsion sheaf supported at finitely
many closed points.  At a point where \(f\) is etale, the relative differential
module vanishes.  Conversely, non-vanishing of \(\Omega_{X/Y}\) measures
failure of etaleness, hence ramification.
\end{proof}

\begin{definition}
The \emph{ramification divisor} of a finite separable morphism of smooth curves
is
\[
    R_f=\sum_{P\in X}\length_{\OO_{X,P}}(\Omega_{X/Y,P})[P].
\]
\end{definition}

\section{Exercises Used Later}

\begin{exercise}[Differentials of a quotient]\label{ex:differentials-quotient}
Let \(B=A[x_1,\ldots,x_n]/(f_1,\ldots,f_r)\).  Show that
\[
    \Omega_{B/A}
    \simeq
    \left(\bigoplus_j B\,dx_j\right)\Big/
    \left(\sum_i B\,df_i\right).
\]
\end{exercise}

\begin{solution}
This is the conormal sequence for the quotient
\[
    A[x_1,\ldots,x_n]\to B.
\]
The polynomial ring has free module of differentials with basis \(dx_j\).  The
ideal \(I=(f_1,\ldots,f_r)\) maps to this free module by
\[
    f_i\longmapsto df_i=\sum_j(\partial f_i/\partial x_j)dx_j.
\]
The cokernel is \(\Omega_{B/A}\).
\end{solution}

\begin{exercise}[Separable field extensions]\label{ex:separable-no-differentials}
Let \(L/K\) be a finite separable field extension.  Show that
\(\Omega_{L/K}=0\).
\end{exercise}

\begin{solution}
It is enough to treat \(L=K(\alpha)\).  Let \(p(T)\) be the minimal polynomial
of \(\alpha\).  Since the extension is separable, \(p'(\alpha)\ne0\).  In
\(\Omega_{L/K}\), the relation \(p(\alpha)=0\) gives
\[
    0=d(p(\alpha))=p'(\alpha)d\alpha.
\]
As \(p'(\alpha)\) is invertible in the field \(L\), \(d\alpha=0\).  Since
\(\alpha\) generates \(L\) over \(K\), all differentials vanish.
\end{solution}

\begin{exercise}[Ramification in a local parameter]\label{ex:ramification-local}
Let \(f:X\to Y\) be a finite morphism of smooth curves, \(P\in X\), and
\(Q=f(P)\).  Choose uniformizers \(t\in\OO_{X,P}\) and \(u\in\OO_{Y,Q}\).  If
\[
    f^\#u=c t^e+\text{higher terms},\qquad c\in\OO_{X,P}^\times,
\]
show that \(e=e_P\).
\end{exercise}

\begin{solution}
The maximal ideal of \(\OO_{Y,Q}\) is generated by \(u\), and its extension to
\(\OO_{X,P}\) is generated by \(f^\#u\).  In the DVR \(\OO_{X,P}\), the order
of \(f^\#u\) is the smallest exponent with non-zero coefficient in its
\(t\)-expansion, namely \(e\).  Therefore
\[
    f^\#(\mathfrak m_Q)\OO_{X,P}=(t^e)=\mathfrak m_P^e,
\]
so \(e=e_P\).
\end{solution}
